% !PRINTED 608
\weberpage{608}
%<<CONT
the roots of the equation
\weq{2}{x^p - a = 0.}

If this equation is reducible, then
\weq{3}{x^p - a = f_1(x) f_2(x),}
and $f_1$, $f_2$ are functions in $\Omega$ of degree lower than $p$.
One part of the roots (1) will fall on $f_1 = 0$, another part on $f_2 = 0$. If
then $\mu$ is the degree of $f_1$, then for some exponent $\lambda$
\[ \varepsilon^\lambda \alpha^\mu = b \]
must be a quantity in $\Omega$ [the term in $f_1(x)$ independent of $x$]; hence, if one
raises to the $p$th power,
\weq{4}{a^\mu = b^p.}

Now $0 < \mu < p$, and therefore $\mu$ and $p$ are relatively prime, so that two integers
$h$ and $k$ can be determined from the equation
\[ \mu h + p k = 1 \]
Then by (4)
\[ a = a^{\mu h} a^{p k} = (b^h a^k)^p; \]
thus, contrary to the hypothesis, $a$ would be the $p$th power of a quantity in $\Omega$.

This theorem, as one sees, is independent of the assumption that $\Omega$ be a
real field. But if we take it to be real, and $p$ is not equal to 2, then the
equation (2) has imaginary roots, and by 1. cannot reduce a normal equation $g(t) = 0$.
But if $p = 2$, then $g(t)$ can be reduced by this equation (2)
only if $a$ is positive and the degree of $g(t)$ is an \gesperrt{even} number; hence:

2.~\gesperrt{A normal equation of odd degree cannot be reduced by the adjunction of a
real radical.}

This case occurs in the casus irreducibilis of the cubic equation, when one has adjoined the
square root of the discriminant to the domain of rationality of the coefficients.
For since the discriminant is positive, the domain of rationality remains
real, and it also remains real however many real radicals one adjoins. If
the equation is thus to be solvable by real radicals, it must at last split under such
adjunctions, which by 2. is impossible.
